% Research continuation at source commit 3447bc59b78d138a7f0d536b66ac6e5cc5d5e49f.
% New result: the global real-order Euler constant is the existing axis maximum.
% This is a self-contained section fragment, not a standalone document.

\section{The sharp universal Euler constant at all positive real orders}
\label{newEuler:sec}

Write
\[
 c_n=\frac{H_{n-1}^{(b)}}{n^a},\qquad
 F_{a,b}(z)=\sum_{n\ge1}c_nz^n,
 \qquad a\ge0,\quad b>0,
\]
where values away from the disk use the principal continuation to
$\mathbb C\setminus[1,\infty)$. Set
\[
 g_{a,b}=\operatorname{Im}F_{a,b}(i),\quad
 d_j=\sum_{k=0}^j(-1)^k\binom jk c_{2k+1},\quad
 E_N=\sum_{j=0}^{N-1}2^{-j-1}d_j\quad(N\ge1).
\]
The manuscript proves one-sided Euler convergence throughout this domain,
but explicitly leaves a universal error constant in the supercritical
fractional-order region open. The Gaussian maximum established there
controls only the first Euler remainder. The following scalar inequality
supplies the missing comparison for every later remainder.

\subsection{An elementary comparison of scalar Euler kernels}

For an integer $N\ge1$ put
\[
 \Phi_N(t)=\frac{(1-t)^N}{1+t}\qquad(0\le t\le1).
\]
\begin{lemma}[Uniform multiplicative-increment bound]
\label{newEuler:lem:scalar}
For every integer $N\ge1$ and $q,t\in[0,1]$,
\begin{equation}\label{newEuler:eq:scalar}
 0\le\Phi_N(qt)-\Phi_N(t)\le\frac{1-q}{1+q}.
\end{equation}
The upper inequality is strict when $0<q<1$ and $0<t<1$.
For $0<q<1$ and $t=1$, equality in the upper inequality occurs
exactly when $N=1$.
\end{lemma}
\begin{proof}
The first inequality follows from the decrease of $\Phi_N$. For the
second define $f_N(u)=u^N/(2-u)$ on $[0,1]$ and set
$A=1-qt$, $B=1-t$, and $p=A+B$. The cases $q=1$ or $t=0$ are
immediate, so suppose $A>B$. The desired inequality is equivalent to
\begin{equation}\label{newEuler:eq:divided}
 \frac{f_N(A)-f_N(B)}{A-B}\le\frac1{2-A-B}.
\end{equation}
The exceptional case $A=B=1$ has already been excluded.

The derivative $f_N'$ is strictly convex on $(0,1)$: expand
\[
 f_N(u)=\sum_{j=0}^{\infty}\frac{u^{N+j}}{2^{j+1}}
\]
and differentiate three times; the result is strictly positive for
$0<u<1$. The average of a convex function over a centered interval
is nondecreasing as its length increases. Explicitly, if
$h\mapsto (2h)^{-1}\int_{c-h}^{c+h}g(v)\,dv$ is differentiated,
the result is nonnegative by the trapezoid bound for a convex $g$.
It is strictly positive for a strictly convex $g$ and a nondegenerate
interval. Applying this to $g=f_N'$ shows that the divided difference
on the left of \eqref{newEuler:eq:divided}, at fixed sum $p$, increases
as the endpoints move apart.

If $p\le1$, enlarge $[B,A]$ to $[0,p]$. Then
\[
 \frac{f_N(A)-f_N(B)}{A-B}
 \le\frac{f_N(p)-f_N(0)}p
 =\frac{p^{N-1}}{2-p}\le\frac1{2-p}.
\]
If $p\ge1$, enlarge it instead to $[p-1,1]$. Since $f_N(1)=1$
and $f_N(p-1)\ge0$,
\[
 \frac{f_N(A)-f_N(B)}{A-B}
 \le\frac{1-f_N(p-1)}{2-p}\le\frac1{2-p}.
\]
When $0<t<1$ and $0<q<1$, both original endpoints are interior
to $[0,1]$, so the interval is strictly enlarged in the appropriate
case and strict convexity gives a strict inequality. At $t=1$ the
original difference equals $(1-q)^N/(1+q)$, proving the equality
statement. The endpoint $q=0$ follows directly from
$\Phi_N(0)-\Phi_N(t)=1-\Phi_N(t)\le1$.
\end{proof}

\subsection{A probability formula for every Euler remainder}

For $a>0$ let $S\sim\operatorname{Gamma}(a,1)$ and
$T\sim\operatorname{Gamma}(b,1)$ be independent, and set
$X=e^{-S}$, $V=e^{-T}$. When $a=0$, define $S=0$, $X=1$
deterministically. A finite geometric sum and the Gamma integral give
\begin{equation}\label{newEuler:eq:moments}
 c_n=\mathbb E\!\left[
       \frac{X^{n-1}-(XV)^{n-1}}{1-V}\right]\quad(n\ge1).
\end{equation}
Indeed the expectation factors as
$n^{-a}\sum_{k=0}^{n-2}(k+1)^{-b}$, including the empty sum at
$n=1$. For $a>0$ this is the existing positive double kernel in
probability coordinates; no signed moment measure is required.

\begin{theorem}[Universal one-sided remainder identity and domination]
\label{newEuler:thm:domination}
For every $a\ge0$, $b>0$ and integer $N\ge1$, the Euler series
converges absolutely, $d_0=0$, $d_j<0$ for $j\ge1$, and
$E_N$ decreases strictly to $g_{a,b}$. More precisely, with
$R_N(a,b)=2^N(E_N-g_{a,b})$,
\begin{equation}\label{newEuler:eq:R}
 R_N(a,b)=\mathbb E\!\left[
 \frac{\Phi_N(X^2V^2)-\Phi_N(X^2)}{1-V}\right]>0.
\end{equation}
Define
\begin{equation}\label{newEuler:eq:C}
 C(b)=\beta(b)+2^{-b}\eta(b)
     =\mathbb E\!\left[\frac{1+V}{1+V^2}\right],
\end{equation}
where $\beta$ is the Dirichlet beta function and
$\eta(b)=\sum_{n\ge1}(-1)^{n-1}n^{-b}$. Then
\begin{equation}\label{newEuler:eq:fixedb}
 R_N(a,b)<C(b)\qquad(a>0,b>0,N\ge1).
\end{equation}
At $a=0$, the same estimate is non-strict, and equality holds exactly
when $N=1$. In addition, $R_N(a,b)\to0$ as $N\to\infty$ for every
fixed pair in the stated domain.
\end{theorem}
\begin{proof}
Applying a finite binomial sum to \eqref{newEuler:eq:moments} gives
\[
 d_j=\mathbb E\!\left[
  \frac{(1-X^2)^j-(1-X^2V^2)^j}{1-V}\right].
\]
This is zero for $j=0$ and strictly negative for $j\ge1$. Every
fixed integrand is bounded by $2j$, by the mean-value theorem, so
all displayed finite expectations are legitimate.

The positive double-kernel continuation at $i$ gives
\begin{equation}\label{newEuler:eq:Gaussian}
 -g_{a,b}=\mathbb E\!\left[
 \frac{X^2(1+V)}{(1+X^2)(1+X^2V^2)}\right]>0.
\end{equation}
This remains valid on the axis $a=0$ from
$F_{0,b}(z)=z\operatorname{Li}_b(z)/(1-z)$.
Alternatively, expanding the two denominators in the disk gives
\eqref{newEuler:eq:moments}, and dominated continuation gives
\eqref{newEuler:eq:Gaussian}.

Tonelli's theorem, applied to the nonnegative functions $-d_j$, permits
summing the infinite Euler series before proving its finiteness. The
geometric sum gives
\[
 \sum_{j\ge0}\frac{-d_j}{2^{j+1}}
 =\mathbb E\!\left[
  \frac{(1+X^2V^2)^{-1}-(1+X^2)^{-1}}{1-V}\right]
 =-g_{a,b}<\infty.
\]
Thus the Euler series is absolutely convergent and its sum is $g$.
The same geometric summation beginning at $j=N$ proves
\eqref{newEuler:eq:R}. The positivity and strict decrease of $E_N$
follow from the strict signs of $d_j$ for $j\ge1$.

Apply Lemma~\ref{newEuler:lem:scalar} with $t=X^2$ and $q=V^2$.
It gives the pointwise domination
\[
 0<\frac{\Phi_N(X^2V^2)-\Phi_N(X^2)}{1-V}
 \le\frac{1-V^2}{(1+V^2)(1-V)}
 =\frac{1+V}{1+V^2}.
\]
The upper inequality is strict almost surely for $a>0$. On the axis
it is equality precisely at $N=1$. Taking expectations proves
\eqref{newEuler:eq:fixedb}. To identify $C(b)$, write the density
of $V$ as $(-\log v)^{b-1}\,dv/\Gamma(b)$ on $(0,1)$, and
integrate the even and odd geometric series for $(1+v)/(1+v^2)$.
Their partial sums are bounded uniformly, justifying dominated
convergence; the resulting alternating Dirichlet series are exactly
$\beta(b)+2^{-b}\eta(b)$ for every $b>0$.

For every $0<X\le1$ and $0<V<1$, both powers in the numerator of
\eqref{newEuler:eq:R} tend to zero as $N\to\infty$.
The same integrable domination, independent of $N$, proves $R_N\to0$.
\end{proof}

\subsection{The least constant uniform in both real orders}

The preceding manuscript already proves that $C$ has exactly one
stationary point $b_*$, a nondegenerate global maximum, with
$1<b_*<2$. Denote its value by
\[
 C_*:=C(b_*)=\max_{b>0}\{\beta(b)+2^{-b}\eta(b)\}.
\]
Its reported numerical diagnostics are
\[
 b_*\approx1.30221658710124120959,\qquad
 C_*\approx1.13656110333950956095;
\]
these decimal values are not needed for the theorem. The existing exact
bounds $227/200<C_*<(1+\sqrt2)/2$ may be retained.

\begin{corollary}[Sharp universal real-order Euler bound]
\label{newEuler:cor:sharp}
For all $a,b>0$ and all integers $N\ge1$,
\begin{equation}\label{newEuler:eq:universal}
 E_N-C_*2^{-N}<g_{a,b}<E_N.
\end{equation}
The number $C_*$ is the least constant for which this assertion holds
uniformly over the entire positive quadrant. With $a=0$ included,
$0<R_N\le C_*$, and equality occurs exactly at
$(a,b,N)=(0,b_*,1)$. For each fixed $b>0$, the least constant uniform
in $a>0$ and $N\ge1$ is $C(b)$.
\end{corollary}
\begin{proof}
The upper bound follows from \eqref{newEuler:eq:fixedb} and
$C(b)\le C_*$. For sharpness, take $N=1$, for which $E_1=0$ and
$R_1=-2g$. As $a\downarrow0$, $S\sim\operatorname{Gamma}(a,1)$
tends to zero in probability (its mean is $a$), so $X\to1$.
The bounded integrand in \eqref{newEuler:eq:Gaussian}, or the
pointwise domination already proved, gives
\[
 \lim_{a\downarrow0}R_1(a,b)=C(b).
\]
Taking $b=b_*$ shows that every smaller universal constant fails for
some positive $a$. The same argument at a fixed $b$ proves its
claimed sharp constant. Strictness and the sole equality case on the
closed outer-order axis follow from the theorem and the uniqueness
of the maximum of $C$.
\end{proof}

The derivative of $(1+v)/(1+v^2)$ is
$(1-2v-v^2)/(1+v^2)^2$, so its maximum on $[0,1]$ is
$(1+\sqrt2)/2<5/4$, attained at $v=\sqrt2-1$.
Theorem~\ref{newEuler:thm:domination} therefore supplies the
self-contained budget $R_N(a,b)<5/4$ throughout the domain.
The explicit algebraic choice $(1+\sqrt2)/2$ gives an immediate
universal enclosure without optimization in $b$.
The established constant one for $a\ge1$ remains sharper on that
restricted region. The new result resolves precisely the outstanding
supercritical fractional-order question; it does not assert that the
scaled errors are monotone in $N$ there, nor that the fixed-pair
constant equals $-2g$ there.

\subsection{A convenient rational certificate for the universal constant}
\begin{proposition}[An elementary rational error budget]
\label{newEuler:prop:rational}
The sharp universal constant satisfies $C_*<57/50$. Consequently
\[
 E_N-\frac{57}{50}\,2^{-N}<g_{a,b}<E_N
 \qquad(a,b>0,\ N\ge1).
\]
\end{proposition}
\begin{proof}
The existing axis-maximum theorem gives $b_*\in(1,2)$. By its positive
Mellin representation, $M(b)=\Gamma(b)C(b)$ is log-convex. Also
\[
 (\log\Gamma)''(b)=\sum_{j\ge0}\frac1{(j+b)^2}
 \le\frac{\pi^2}{6}<\frac53\qquad(1\le b\le2).
\]
On an interval $[u,u+h]\subset[1,2]$, the error of linear
interpolation of $\log\Gamma$ is therefore at most $5h^2/24$.
Combining this with the convexity inequality for $\log M$ yields
\[
 C(b)\le\max\{C(u),C(u+h)\}\exp(5h^2/24).
\]
The accompanying exact standard-library verifier
\texttt{certify\_rational\_euler\_constant.py} proves
\[
 C(k/10)<\frac{1137}{1000}\qquad(k=10,11,\ldots,20).
\]
It uses 80 terms of the axis Euler enclosure with error
budget $5/4$, which follows directly from the scalar domination above
and was also available in the preceding manuscript. Every $n^{-k/10}$ is enclosed on a $2^{-160}$ grid
by integer inequalities: if
$r=\lfloor (2^{1600}/n^k)^{1/10}\rfloor$, then
$r^{10}n^k\le2^{1600}<(r+1)^{10}n^k$. All subsequent operations
are exact rational arithmetic. The full rational intervals are saved
in \texttt{rational\_euler\_constant\_certificate.json}; none of the
displayed decimal diagnostics is used for a sign decision.
Taking $h=1/10$ gives
\[
 C_*<\frac{1137}{1000}e^{1/480}
 <\frac{1137}{1000}\frac{480}{479}
 =\frac{13644}{11975}<\frac{57}{50}.
\]
The exponential inequality follows from its positive series and
$e^x<(1-x)^{-1}$ for $0<x<1$.
\end{proof}

\begin{proposition}[A narrow exact enclosure of the sharp constant]
\label{newEuler:prop:narrow}
The unique maximizer and its value satisfy
\[
 1.3022165871012412095923717014
 <b_*<1.3022165871012412095923717017,
\]
\[
 \begin{gathered}
 1.136561103339509560952586375779942387681723540765044951\\
 <C_*<\\
 1.136561103339509560952586375779942387681723540765044952.
 \end{gathered}
\]
These are exact rational enclosures, not estimates from a
floating-point stationary-point calculation.
\end{proposition}
\begin{proof}
Let $L,R$ be the displayed endpoints for $b_*$, and put
$c=1.30221658710124120959237170155$. The standard-library verifier
\texttt{certify\_axis\_maximum.py} encloses $C(L)$, $C(c)$ and
$C(R)$ by exact rational intervals and proves
\[
 C(c)>C(L),\qquad C(c)>C(R).
\]
The previously proved strict unimodality implies $L<b_*<R$.

Here are all arithmetic ingredients of that finite certificate.
The computation uses 224 axis Euler terms, with the self-contained
error budget $5/4$ proved above. To enclose $n^{-b}$ for rational $b$, write
$n=2^k r$, $1\le r<2$, and use
\[
 \log n=k\log2+2\sum_{j=0}^{127}\frac{z^{2j+1}}{2j+1}
             +\mathcal R,\qquad z=\frac{r-1}{r+1},
 \quad 0\le\mathcal R\le\frac{2z^{257}}{257(1-z^2)}.
\]
The same formula with $z=1/3$ encloses $\log2$. For each resulting
$0\le x=b\log n<8$, sum the positive exponential series through
$x^{200}/200!$, with tail bounded by
\[
 \frac{x^{201}}{201!}\frac1{1-x/202},
\]
and reciprocate the interval. Every operation is rounded outwards
on the dyadic grid $2^{-384}$. All sum weights and error budgets
are exact rational numbers. The complete rational intervals, along
with the strict comparisons above, are recorded in
\texttt{axis\_maximum\_certificate.json}.

For the upper enclosure on the maximum, the interpolation argument
in Proposition~\ref{newEuler:prop:rational} gives, with
$\delta=5(R-L)^2/24$,
\[
 C(c)\le C_*\le\max\{C(L),C(R)\}e^{\delta}
 <\frac{\max\{C(L),C(R)\}}{1-\delta}.
\]
Applying the rational lower endpoint for $C(c)$ and the two rational upper
endpoints at $L,R$, then rounding the final bounds outwards to
54 decimal places, proves the displayed enclosure.
\end{proof}
